% !TeX program = xelatex
% ============================================================
%  2. ЦЕЛЬ И ЗАДАЧИ РАБОТЫ
% ============================================================

\section{Цель и задачи работы}
\label{sec:goal}

\subsection{Цель}

\begin{center}
\fbox{\parbox{0.94\textwidth}{%
\textbf{Цель:} разработать численный метод оценки бахадуровской эффективности критериев
согласия, реализовать его в виде нового типа эксперимента фреймворка
\texttt{pysatl-experiment} и подтвердить корректность метода сопоставлением численных
оценок с аналитическими значениями для критериев, для которых такие значения получены.
}}
\end{center}

\subsection{Декомпозиция работы на задачи}

\begin{enumerate}[label=\textbf{\arabic*.},leftmargin=1.9cm]

  \item \textbf{Подготовительная.} \planned{5 семестр, спринт 0}
    \begin{enumerate}[label=\arabic*.\arabic*,leftmargin=1.9cm]
      \item Согласовать тему, цель и план работы с руководителем.
      \item Договориться с коллегой по команде, кто что делает: аналитические значения
            эффективности считает она, численный метод и его встраивание во фреймворк~---
            я; в каком виде она передаёт мне свои значения для сравнения.
      \item Решить с командой, в каком из двух репозиториев будет находиться код метода
            и как он попадёт к пользователям (\texttt{pysatl-criterion} подключается как
            отдельный пакет, поэтому его изменение требует публикации новой версии).
      \item Завести репозиторий отчётности, настроить автоматическую сборку PDF, написать
            аннотацию, цель и задачи.
    \end{enumerate}

  \item \textbf{Разобраться в методе и выбрать способ расчёта.} \planned{5 семестр, спринт 1--2}
    \begin{enumerate}[label=\arabic*.\arabic*,leftmargin=1.9cm]
      \item Изучить, что такое бахадуровская эффективность и как она считается;
            разобраться, почему для большинства критериев библиотеки классическая формула
            наклона не применима, и выбрать рабочий вариант~--- приближённый наклон.
      \item Определить, какие две величины нужны для расчёта и как получить каждую из них.
      \item Вывести, как из этих величин складываются наклон и эффективность~---
            относительная (критерий против критерия) и абсолютная (критерий против
            наилучшего возможного теста).
    \end{enumerate}

  \item \textbf{Доработка фреймворка под новую задачу.} \planned{5 семестр, спринт 3--4}
        Есть существующие ограничения, которые мешают реализации нового типа эксперимента; их нужно устранить.

  \item \textbf{Реализация нового типа эксперимента.} \planned{5 семестр, спринт 5--6}
    \begin{enumerate}[label=\arabic*.\arabic*,leftmargin=1.9cm]
      \item Настроить вход: новое значение типа эксперимента, его описание в конфигурации
            и проверка того, что необходимые данные уже посчитаны.
      \item Реализовать расчёт: воркер, который считает наклон и эффективность, и его
            запуск в несколько потоков.
      \item Реализовать сохранение результатов в базу и повторное использование уже
            посчитанного при перезапуске.
      \item Собрать всё вместе, чтобы эксперимент создавался и запускался из командной
            строки как существующие.
      \item Сделать отчёт: таблицу эффективностей, график зависимости от объёма выборки,
            сравнение с аналитическими значениями.
      \item Добавить пример конфигурации, документацию и тесты.
    \end{enumerate}

  \item \textbf{Проверка метода и эксперименты.} \planned{6 семестр}
    \begin{enumerate}[label=\arabic*.\arabic*,leftmargin=1.9cm]
      \item Сравнить численные оценки с аналитическими значениями коллеги; если есть
            расхождения~--- разобраться в их причинах.
      \item Проверить, что результат не меняется при увеличении числа итераций и объёма
            выборки; оценить, сколько времени занимает расчёт.
      \item Провести итоговый расчёт для выбранного семейства распределений и набора
            альтернатив, объяснить получившееся ранжирование критериев.
    \end{enumerate}

  \item \textbf{Оформление результатов и подготовка к защите.} \planned{6 семестр}
    \end{enumerate}

\end{enumerate}

\paragraph{Что не входит в задачи.} Обзор литературы отдельной задачей не является: он
ведётся непрерывно начиная со спринта~2 и оформляется файлом \texttt{03-materials.tex}.
Аналитический вывод формул и расчёт эталонных значений эффективности выполняет коллега
по команде.
